\subsection{有限高斯波包：解析解 TDSE 与时间聚焦}
\label{sec:electron-gaussian-lens-analytic}

时间透镜的聚焦也可以直接从一个有限波包的波函数计算。取已建立的纵向二阶色散，在随中心电子运动的坐标 $\zeta=z-v_0t$ 中，包络满足含时 Schrödinger 方程（time-dependent Schrödinger equation, TDSE）
\begin{equation}
 \ii\hbar\partial_t\psi(\zeta,t)
 =-\frac{\hbar^2}{2m_\parallel}\partial_\zeta^2\psi(\zeta,t),
 \qquad \int|\psi|^2\dd\zeta=1.
 \label{eq:electron-gaussian-envelope-tdse}
\end{equation}
这个方程与自由非相对论粒子的 Schrödinger 方程形式相同，但系数 $m_\parallel=\gamma_0^3m_{\mathrm e}$ 来自相对论能量的二次展开。中心能量和匀速运动已被移入载波相位，$\psi$ 只描述相对于中心电子的包络变化。它的适用条件仍是实际动量范围内高阶色散积累的相位可以忽略。

用复参数 $a_0$ 表示一般纯高斯输入，$\Re a_0>0$：
\begin{equation*}
 \psi(\zeta,0)=\mathcal N_0\ee^{-a_0\zeta^2},\qquad
 \mathcal N_0=\left(\frac{2\Re a_0}{\pi}\right)^{1/4}.
\end{equation*}
$a_0$ 的实部决定密度宽度，虚部决定二次空间相位；$[a_0]=\mathrm{m^{-2}}$。连续 TDSE 的精确高斯解为
\begin{equation}
 \psi(\zeta,t)=\frac{\mathcal N_0}{\sqrt{d(t)}}
 \exp\!\left[-\frac{a_0\zeta^2}{d(t)}\right],\qquad
 d(t)=1+\frac{2\ii\hbar a_0t}{m_\parallel}.
 \label{eq:electron-gaussian-tdse-solution}
\end{equation}
平方根支路从 $d(0)=1$ 连续选取。$d$ 无量纲；前因子同时包含保持归一化所需的幅度变化与共同传播相位，不能只保留指数中的宽度变化。

\begin{exbox}[带二次相位的波包何时最窄]
令初始密度标准差为 $\sigma_{\zeta0}$，空间啁啾参数为实数 $C_\zeta$，即
\begin{equation*}
 a_0=\frac{1-\ii C_\zeta}{4\sigma_{\zeta0}^2},\qquad
 t_{\rm dif}=\frac{2m_\parallel\sigma_{\zeta0}^2}{\hbar},\qquad
 u_t=\frac{t}{t_{\rm dif}}.
\end{equation*}
这里 $C_\zeta$ 是无量纲实数，控制位置表象中的二次相位；前面用 $C_p$ 表示的是动量表象中的二次相位。$t_{\rm dif}$ 具有时间单位，它是无啁啾高斯包络因自由传播而明显展宽的时间尺度，$u_t$ 则以这个时间为单位。代入一般解，有 $d=1+(C_\zeta+\ii)u_t$。利用
\begin{equation*}
 \Re\frac{a_0}{d}=\frac{\Re a_0}{|d|^2},\qquad
 |d|^2=(1+C_\zeta u_t)^2+u_t^2,
\end{equation*}
便从概率密度的高斯指数读出
\begin{equation*}
 \sigma_\zeta^2(t)=\sigma_{\zeta0}^2
 \left[1+2C_\zeta u_t+(1+C_\zeta^2)u_t^2\right].
\end{equation*}
这是关于时间的二次式，配方即可求极小值：
\begin{equation*}
 \frac{\sigma_\zeta^2(t)}{\sigma_{\zeta0}^2}
 =(1+C_\zeta^2)
 \left(u_t+\frac{C_\zeta}{1+C_\zeta^2}\right)^2
 +\frac1{1+C_\zeta^2}.
\end{equation*}
若 $C_\zeta<0$，未来存在焦点，
\begin{equation}
 t_f=-\frac{C_\zeta}{1+C_\zeta^2}t_{\rm dif},\qquad
 L_f=v_0t_f,\qquad
 \sigma_{\zeta,\min}=\frac{\sigma_{\zeta0}}{\sqrt{1+C_\zeta^2}}.
 \label{eq:electron-gaussian-focus}
\end{equation}
若 $C_\zeta=0$，波包从初始时刻起展宽；若 $C_\zeta>0$，二次式的焦点位于过去，未来同样单调展宽。因此聚焦需要正确符号的相位相关，而不只是足够大的能宽。

若输入本来无啁啾，在时间透镜的相位极值附近乘入 $\exp[-\ii|\beta|\omega^2\tau^2]$，并使用 $\zeta=-v_0\tau$，便得到
\begin{equation*}
 C_\zeta=-4|\beta|\omega^2\sigma_{\tau0}^2,
 \qquad \sigma_{\tau0}=\sigma_{\zeta0}/v_0.
\end{equation*}
其符号恰好允许未来聚焦；将它代入\eqnrefs{eq:electron-gaussian-focus} 就得到对任意允许参数成立的焦距与压缩比。到达时间宽度取 $\sigma_\tau\simeq\sigma_\zeta/v_0$，还要求速度展宽相对中心速度很小。
\end{exbox}

为了理解压缩过程中动量分布怎样变化，再计算位置与动量的协方差。记 $\Delta\zeta=\hat\zeta-\langle\hat\zeta\rangle$、$\Delta p=\hat p-\langle\hat p\rangle$；由于两个算符不对易，用对称乘积定义实协方差。对上述初态，在位置表象使用 $\hat p=-\ii\hbar\partial_\zeta$，得
\begin{equation*}
 \sigma_p^2=\frac{\hbar^2(1+C_\zeta^2)}{4\sigma_{\zeta0}^2},\qquad
 C_{\zeta p}\equiv\tfrac12\langle\Delta\zeta\Delta p+\Delta p\Delta\zeta\rangle
 =\frac{\hbar C_\zeta}{2}.
\end{equation*}
所以 $\sigma_{\zeta0}^2\sigma_p^2-C_{\zeta p}^2=\hbar^2/4$。自由传播保持 $\sigma_p$ 不变，只把协方差变为 $C_{\zeta p}(t)=C_{\zeta p}(0)+t\sigma_p^2/m_\parallel$；焦点恰在该协方差过零时，且 $\sigma_{\zeta,\min}\sigma_p=\hbar/2$。因此焦点处位置与动量不再相关，宽度乘积恰好达到最小不确定值。

时间透镜的适用条件也应作用于\emph{相位本身}。若检查包络的 $|\tau|\le K\sigma_{\tau0}$ 区域，其中 $K$ 是选定的覆盖倍数，则正弦相位余项的首项为
\begin{equation*}
 |\delta\varphi_4|\lesssim\frac{|\beta|}{12}
 (\omega K\sigma_{\tau0})^4.
\end{equation*}
要可靠地用二次相位替代完整作用，应同时检查 $\omega K\sigma_{\tau0}\ll1$ 和 $|\delta\varphi_4|\ll1$，并检查高阶色散在 $t_f$ 内的累积误差。仅有局部几何孔径很小，在极强驱动下仍不足以保证波函数误差很小。

\begin{derivation}[从自由传播核到高斯解，并直接代回 TDSE]
在二阶包络模型中，连续动量谱给传播核
\begin{equation*}
 K(\zeta,\zeta';t)=\int\frac{\dd p}{2\pi\hbar}
 \exp\!\left[\frac{\ii p(\zeta-\zeta')}{\hbar}
             -\frac{\ii p^2t}{2m_\parallel\hbar}\right]
 =\sqrt{\frac{m_\parallel}{2\pi\ii\hbar t}}
 \exp\!\left[\frac{\ii m_\parallel(\zeta-\zeta')^2}{2\hbar t}\right],
 \quad t>0.
\end{equation*}
振荡高斯积分可先加正的收敛因子，再取其趋零极限。将初态代入 $\psi(\zeta,t)=\int K\psi(\zeta',0)\dd\zeta'$，整理为
\begin{equation*}
 -A\zeta'^2+B\zeta'+C,
 \quad A=a_0-\frac{\ii m_\parallel}{2\hbar t},\quad
 B=-\frac{\ii m_\parallel\zeta}{\hbar t},\quad
 C=\frac{\ii m_\parallel\zeta^2}{2\hbar t}.
\end{equation*}
$\Re A>0$ 保证积分收敛；使用 $\int\ee^{-Ax^2+Bx}\dd x=\sqrt{\pi/A}\ee^{B^2/(4A)}$，合并前因子与指数即得\eqnrefs{eq:electron-gaussian-tdse-solution}。

还可独立代回检验。写 $a(t)=a_0/d(t)$、$\mathcal N(t)=\mathcal N_0/\sqrt{d(t)}$，则
\begin{equation*}
 \dot a=-\frac{2\ii\hbar}{m_\parallel}a^2,\qquad
 \frac{\dot{\mathcal N}}{\mathcal N}=-\frac{\ii\hbar}{m_\parallel}a,
 \qquad
 \frac{\partial_\zeta^2\psi}{\psi}=-2a+4a^2\zeta^2.
\end{equation*}
TDSE 左右两侧的常数项与 $\zeta^2$ 项于是分别相等。再用
$\Re(a_0/d)=\Re a_0/|d|^2$，对 $|\psi|^2$ 作高斯积分，归一化保持为一。所以传播后的函数同时满足初始条件、运动方程和归一化条件。
\end{derivation}
